% Ampliación argumental. Insertar tras la introducción de excepcional.tex,
% antes de su primera subsección. Conserva íntegro el desarrollo existente.
\paragraph{Articulation des coordonnées régionales et de l'incidence de clôture.}
\label{inc-ampl-articulacion}
La fonction de ce prolongement se comprend à partir d'une distinction
entre trois opérations : conserver un registre, évaluer une coordonnée et
extraire une relation d'incidence. Toutes trois agissent sur une information
produite par APP--TRIT--TPK. Leur articulation explicite la thèse
structurelle de l'article : une réalisation numérique possède une provenance
opératorielle, et certaines données de cette provenance admettent en outre une
réalisation combinatoire et réticulaire. Les équations permettent de suivre les deux
parcours avec leurs domaines et avec l'information qui subsiste dans chacun.

Le point de départ est l'état terminal muni de ses prolongements
compatibles. Il conserve les restes et quotients des deux feuillets
APP, le régime et l'orientation TRIT, ainsi que le transport TPK avec sa mémoire.
Les régions de clôture, de propagation et d'auto-échelle se déterminent dans ce
domaine avant que leurs lecteurs établissent les coordonnées
$\pi,e,\varphi$. Simultanément, le registre orienté des douze
fenêtres détermine le vecteur signé $U$. L'apparition ultérieure de $K$
provient de la transformation entière démontrée dans
\S\ref{subsec:k-hadamard} :
\begin{equation}
 U=\mathcal H_{12}K,\qquad
 K=\frac14\mathcal H_{12}U,\qquad
 \mathcal H_{12}^{2}=4I_{12}.
 \label{inc-ampl-registro}
\end{equation}
La congruence qui définit $\mathcal L_H$ garantit l'intégralité de
l'inverse. Ainsi, $K$ constitue une mise en coordonnées réversible
du registre signé dans le sous-réseau pertinent. Ses positions conservent
l'ordre de phase ; ses différences $D_3K,D_4K$ conservent les relations
transversales ; la somme $Q(K)$ détermine la composante uniforme. Cette
structure explique pourquoi le registre intervient à la fois dans la composition
numérique et dans la sélection d'une configuration de frontière.

L'évaluation périodique
\[
 K\longmapsto\kappa_{\mathrm{per}}
 =\frac{\sum_{m=1}^{12}K_m1000^{12-m}}{1000^{12}-1}
\]
maintient les douze blocs récupérables lorsque la base, la longueur et
l'origine de phase sont fixées. L'extraction d'un support réalise une opération
différente : elle conserve l'appartenance de chaque position à une classe définie
par un seuil et regroupe les registres qui possèdent le même support. La
distinction entre ces deux applications est décisive. La précision
scalaire de $\kappa_{\mathrm{per}}$ et l'incidence de $B_K$ sont deux
contenus mathématiques spécifiques, reliés par le registre
complet dont ils proviennent.

\paragraph{La normalisation d'alpha et sa configuration signée.}
\label{inc-ampl-clausura}
La coordonnée de $\alpha$ est déterminée dans la voie positionnelle par la
composition des blocs régionaux et du registre dodécaphasique. Avant
la division positionnelle, la combinaison orientée possède les composantes
$s_j=P_j+E_j-\Phi_j-K_j^{\mathrm{per}}$. La normalisation vérifie
\begin{equation}
 A_j=s_j+c_{j+1}-1000c_j,\qquad 0\leq A_j<1000.
 \label{inc-ampl-normalizacion}
\end{equation}
Cette identité conserve simultanément le bloc normalisé et la
contribution du quotient transporté. La limite des préfixes
compatibles détermine la coordonnée scalaire de la clôture. La
configuration signée conserve, quant à elle, les positions de la
combinaison antérieure à la normalisation qui présentent un défaut négatif.
Dans la fenêtre initiale, cette configuration est le vecteur $s_0$ de
\eqref{exc:precarry}, et son support négatif est $H^-_\alpha$.

Deux contenus d'une même opération sont ainsi définis. La
normalisation positionnelle détermine une coordonnée réelle au moyen de ses
restes et quotients compatibles ; le support négatif détermine une
configuration de positions dans le cycle de douze composantes. Le signe
se rapporte au bilan des registres antérieur à la normalisation. Sa lecture
incidencielle reste disponible parce que l'état conserve ce bilan
avec le résultat normalisé. Cette articulation correspond à la
voie positionnelle démontrée ; la comparaison avec les opérateurs analytiques
de lacunes conserve la portée précise établie dans le chapitre
consacré à $\alpha$.

\paragraph{Du codage ternaire au drapeau marqué.}
\label{inc-ampl-bandera}
Les régions fournissent également les codes $w_\pi,w_e,w_\varphi$ et
$w_{\rm APP}$ dans le repère hérité. L'application
\[
 \gamma_{A_W}:w\longmapsto(w,wA_W)
\]
conserve le mot de six composantes comme première coordonnée et ajoute
sa transformation par $A_W$. La relation $A_W^2=-I_6$ fixe la
compatibilité quadratique de cette réalisation. Le pas suivant,
$c\mapsto\operatorname{supp}(c)$, conserve les positions non nulles.
Pour les mots de poids six, le support identifie exactement la paire
$\{c,-c\}$, comme le prouve la distance minimale du code. Ainsi,
le quotient par le signe produit les hexades, tandis que le code complet
reste disponible pour les opérations qui requièrent ses symboles.

Dans cette coordinatisation, les deux constructions de la tétrade coïncident :
\begin{equation}
 \{i:K_i\geq729\}
 =H_e\cap P_\pi=B_K=\{4,6,9,10\}.
 \label{inc-ampl-cuaterna}
\end{equation}
Le membre de gauche provient d'une comparaison entière des composantes
de $K$ ; celui de droite, de l'intersection des supports codés de
propagation et de clôture. Le seuil appartient à la coordinatisation par blocs
déclarée et reste explicité dans la définition du lecteur. Le
registre intégral conserve les valeurs sur lesquelles il agit, même lorsque
des seuils distincts produisent le même sous-ensemble. L'égalité exprime
donc une compatibilité entre des applications spécifiées.

Le support $H^-_\alpha$ complète cette tétrade en une hexade. La
condition d'appartenance au support $P_\pi$, avec les tailles
d'intersection $(4,3,2)$ relativement à $H_e,H_\varphi,H_{\rm APP}$,
sélectionne une unique hexade. Les preuves suivantes réunissent les deux
déterminations : celle obtenue par le signe de $s_0$ et celle obtenue
par l'incidence. Le résultat s'organise sous la forme
\begin{equation}
 \widetilde{\mathcal I}_*
 =\bigl(B_K\subset H^-_\alpha;
         P_\pi,H_e,H_\varphi,H_{\rm APP}\bigr).
 \label{inc-ampl-marco}
\end{equation}
L'inclusion est le drapeau ; les quatre supports supplémentaires constituent
son repère marqué. Les permutations simultanées transportent le repère et
conservent ses inclusions et ses intersections. L'information pertinente réside
dans cette configuration transportée et dans sa classe d'isomorphisme.

Les cinq régions de $\pi$ retrouvent ici leur différenciation interne.
Elles partagent $w_\pi$ comme code ternaire visible, mais conservent des émissions,
des signatures et des multiplicités distinctes. Dans la réalisation binaire sur une
dodécade marquée, la région transversale correspond à l'octade dont
l'intersection avec la dodécade est la tétrade ; les quatre autres régions
correspondent aux relèvements des quatre hexades de son étoile. La
région négative est distinguée par $H^-_\alpha$ et l'ordre de coordinatisation
individualise les trois régions positives. Cette correspondance conserve
les rôles régionaux et la distribution $432+4\cdot144=1008$, tout en
réalisant la partition d'incidence $1+4$.

\paragraph{De l'origine incidencielle à la modification réticulaire.}
\label{inc-ampl-reticulo}
Le drapeau marqué détermine une position au moyen d'une opération
ensembliste équivariante :
\[
 o_*=(H_e\cap H_\varphi)
       \setminus(H^-_\alpha\cup H_{\rm APP})=\{1\}.
\]
Sa fonction consiste à distinguer une composante parmi les douze qui
interviennent dans la réalisation réticulaire. Le code ternaire complet
agit comme code de recollement dans le module discriminant
$(A_2^*/A_2)^{12}$. Sa préimage définit $N(C_W)$, et
l'autodualité, les poids et la distance du code se traduisent,
respectivement, en propriétés d'intégralité et de déterminant, de parité
et en bornes de norme des classes de recollement. L'information symbolique
du code remplit ici une fonction que le support isolé ne remplit plus.

La métrique $A_2$ et la chambre radiale positive sont fixées dans le développement
réticulaire. À l'intérieur de cette chambre, la congruence nécessaire pour un voisin
pair d'indice trois possède douze réalisations de norme minimale $54$,
permutées par le changement de composante distinguée. L'origine $o_*$
sélectionne
\[
 v_\alpha=4\rho_1+\rho_2+\cdots+\rho_{12},\qquad
 v_\alpha^2=2(4^2+11)=54.
\]
Le coefficient $4$ appartient à cette solution de la congruence dans la
chambre déclarée. L'indice $\alpha$ enregistre la provenance
incidencielle de la sélection : son contenu est le repère qui est intervenu dans
la clôture et qui individualise à présent une direction réticulaire.

L'appariement avec $v_\alpha$ détermine le sous-réseau
$N_{\alpha,0}$ d'indice trois, et l'adjonction de la classe
$v_\alpha/3$ produit le voisin $N_\alpha$ de
\eqref{exc:vecino}. Cette modification conserve le rang et rétablit
l'unimodularité et la parité ; la sélection par congruence élimine les racines
antérieures, et les bornes locales excluent les racines dans les nouvelles classes.
L'identification à Leech intervient après ces vérifications. La
réalisation binaire des cinq régions et cette construction réticulaire
sont deux prolongements du repère commun : la seconde reçoit directement
le code ternaire comme donnée de recollement.

\paragraph{Graduation, unité et automorphismes.}
\label{inc-ampl-automorfismos}
La construction de Frenkel--Lepowsky--Meurman s'applique au réseau
$N_\alpha$ une fois ses propriétés vérifiées. Elle incorpore l'espace
des oscillateurs, l'extension centrale du réseau, le produit de vertex
et le secteur tordu par le relèvement de la négation. Le résultat est
l'algèbre $V^\natural_{\rm HMT}$ de \eqref{exc:flm}, dont le groupe
d'automorphismes est isomorphe au Monstre. Son domaine d'action est ainsi
fixé par une structure algébrique munie d'une graduation et d'un produit. Le théorème
de Moonshine de Borcherds détermine ensuite les traces graduées
avec insertion de ces automorphismes \cite{flm,borcherds}.

L'unité acquiert à ce niveau une réalisation précise : la composante
de poids zéro est la droite du vide et fournit le coefficient $1$ de
$q^{-1}$ dans le caractère gradué. Au niveau cylindrique, les
applications de raffinement conservent la fonction constante par
précomposition. Chaque assertion de conservation possède donc son
objet et son application. Le développement des preuves rend visible cette
continuité de construction, depuis les opérations arithmétiques et la
mémoire jusqu'à l'incidence, à la métrique et aux produits de vertex.

Cette articulation situe $K$ et la clôture de $\alpha$ dans l'argument
principal. $K$ conserve de manière réversible une coordonnée entière de
l'histoire ; la normalisation de $\alpha$ compose les registres régionaux
avec cette coordonnée ; la configuration signée détermine un drapeau ; et
le drapeau marqué sélectionne la modification réticulaire sur laquelle se
réalise Moonshine. Les démonstrations suivantes développent ces
applications, en distinguant les équivalences réversibles, les quotients
d'information et les réalisations qui ajoutent une métrique ou une structure
algébrique. Ainsi s'exprime une relation démontrable entre structure,
équation et valeur, en conservant l'attribution de chaque construction et la
portée de ses théorèmes.
